% Generated draft: requires mathematical and source review.
\documentclass[11pt]{article}
\usepackage[margin=25mm]{geometry}
\usepackage{fontspec}
\setmainfont{Latin Modern Roman}
\usepackage{amsmath,amssymb,mathtools,unicode-math}
\setmathfont{Latin Modern Math}
\usepackage{enumitem,needspace,xurl,hyperref}
\setlength{\parindent}{0pt}
\setlength{\parskip}{5pt}
\setlength{\emergencystretch}{3em}
\newcommand{\N}{\mathbb{N}}
\newcommand{\Z}{\mathbb{Z}}
\newcommand{\Q}{\mathbb{Q}}
\newcommand{\R}{\mathbb{R}}
\newcommand{\C}{\mathbb{C}}
\newcommand{\F}{\mathbb{F}}
\newcommand{\A}{\mathbb{A}}
\newcommand{\PP}{\mathbb{P}}
\newcommand{\E}{\mathbb{E}}
\newcommand{\eps}{\varepsilon}
\newcommand{\dd}{\,\mathrm{d}}
\begin{document}
\section*{UnsolvedMath ID: 20002697}
\textbf{Source problem code:} catalogue entry
\subsection{Definitions and mathematical statement}
The source asks for structural explanations of Coxeter--Catalan coincidences among noncrossing partitions $NC_W$, nonnesting objects $NN_W$, and facets of the cluster complex $\Delta_W$. In a finite crystallographic Coxeter setting, $NC_W$ is the interval from $1$ to a Coxeter element in absolute/reflection order, $NN_W$ consists of antichains of the positive-root poset, and $\Delta_W$ is the cluster complex on almost-positive roots. The noncrystallographic meaning of the root-poset model is not specified in the extract. The source states common Catalan cardinalities and common Narayana rank/height/$h$-numbers, asks for statistic-preserving bijections and a general algebraic framework, asks how to formalize Bessis's ``dual'' Coxeter systems, and asks for an invariant-theoretic meaning of $\operatorname{Cat}(W)$, noting that exponents equal invariant degrees minus one. Its no-known-formula/no-theoretical-connection language is historical source wording, not a current assessment.

The remarks add two refinements. Chapoton's $M$, $F$, and $H$ triangles should satisfy explanatory structural relationships. Also, for the Coxeter intersection lattice $\Pi_W$ and maps
\[
 f(w)=\operatorname{Fix}(w),\qquad
 g(A)=\bigcap_{\alpha\in A}\alpha^\perp,
\]
the source reports equinumerous inverse images of every $W$-orbit in $\Pi_W$. It asks for a theoretical proof, an analogous cluster statistic, and compatibility with the triangle refinements. It describes an explicit Reading bijection from $NC_W$ to cluster facets with a then-case-by-case proof; that remark is truncated before its final reference context.

The aid supplies only the Boolean special case $W=A_1^n$. Take coordinate roots $\alpha_i$ and commuting coordinate reflections $s_i$. For each $I\subseteq\{1,\ldots,n\}$ define
\[
 w_I=\prod_{i\in I}s_i,\qquad
 A_I=\{\alpha_i:i\in I\},\qquad
 C_I=\{\alpha_i:i\in I\}\cup\{-\alpha_i:i\notin I\}.
\]
These index all noncrossing elements, root antichains, and cluster facets, respectively. The reported correspondence preserves reflection rank and antichain size $|I|$. Order the facets by any linear extension of subset inclusion. The shelling restriction face of $C_I$ is its positive-root part, so its restriction statistic is also $|I|$. Defining $q(C_I)$ as the intersection of the hyperplanes perpendicular to that positive part gives the exact-flat identity
\[
 \operatorname{Fix}(w_I)=g(A_I)=q(C_I)
 =\{x:x_i=0\text{ for }i\in I\}.
\]
The common multivariate flat enumerator is $\sum_I\prod_{i\in I}z_i=\prod_i(1+z_i)$; the common Narayana/$h$-polynomial is $(1+t)^n$, and the common total is $2^n$.

With $F(x,y)$ counting faces by their positive and negative roots, $H(x,y)=\sum_Ax^{|A|}y^{|A\cap\{\alpha_i\}|}$, and
\[
 M(x,y)=\sum_{I\subseteq J}(-1)^{|J\setminus I|}x^{|J|}y^{|I|},
\]
the aid's convention-explicit Boolean factorizations are
\[
 F=(1+x+y)^n,\qquad H=(1+xy)^n,\qquad M=(1-x+xy)^n.
\]
They do not identify the triangles as equal functions. This exact-flat Boolean example does not give the requested general-type bijections, uniform refined proof, dual-Coxeter framework, or invariant-theoretic interpretation; the full general Chapoton relations are not stated in the supplied record.
\subsection{Short English statement}
The broad question seeks structural and refined explanations of Coxeter--Catalan numerology. The aid treats only products of rank-one systems, where subsets give exact-flat bijections and explicit triangle factorizations. The source's general algebraic, invariant-theoretic, and refinement questions remain separate.
\subsection{Sources}
\begin{itemize}
\item Local UnsolvedMath registry, permanent ID 20002697, source code catalogue entry; created and updated 14 August 2026. Canonical \texttt{aim-probability-notes.json}, zero-based \texttt{notes[138]}, tag \texttt{problem}, item 1.1, ``Explain the numerology.''
\item AIM ``Braid groups, clusters and free probability'': \url{https://aimath.org/WWN/braidgroups/braidgroups.pdf}. The supplied verification reports restoring $NC_W$, $NN_W$, $\Delta_W$, and the map $g$ from the official PDF; statement status is reconstructed\_unverified. No external inspection is claimed here.
\item Original problem/remarks cite David Bessis [6], invariant theory [31], Fr\'ed\'eric Chapoton [17,18], Christos Athanasiadis--Vic Reiner [4], and Nathan Reading [41]. Their full bibliography is absent; the Reading remark ends mid-parenthesis. The source says the orbit refinement was proved case-by-case with computer assistance in exceptional types.
\item Research attempt 1: \texttt{partial\_result}, problem status at run \texttt{partially\_solved}, rigor label \texttt{proved}, novelty confidence low. Ulam AI / UnsolvedMath Contributors, CC BY 4.0; machine-generated aid, not peer reviewed, independently verify before citation.
\end{itemize}
\end{document}
